% Chapter 3 -- Rubric: PO2 (2.1), PO3 (3.1, 3.2)
\chapter{Requirements Analysis}
\label{ch:requirements}

\section{Project Objectives and Constraints}
\label{sec:requirements}

\subsection{Stakeholders and their expectations}

\begin{table}[htbp]
  \centering
  \caption{Stakeholders, expectations, and conflicts}
  \label{tab:stakeholders}
  \begin{tabularx}{\textwidth}{>{\hsize=0.55\hsize}Y>{\hsize=1.2\hsize}Y>{\hsize=1.25\hsize}Y}
    \toprule
    Stakeholder & Expectation & Conflict with other stakeholders \\
    \midrule
    Admission candidate (primary user) & Shortest path to a higher score; immediate practice; questions that feel like the real examination; no wasted time. & Wants to start practising at once; the system wants a 30-question diagnostic first. Wants encouragement; the most informative question is the one they are about 50--70\% likely to fail. \\
    Teacher / coaching instructor & Syllabus fidelity, factually correct content, results that map to chapters they teach. & Thinks in chapters; the engine thinks in skills that cut across chapters. Requires content correctness guarantees that a generative pipeline cannot fully give. \\
    Project team (developer) & A system that is demonstrably correct, maintainable, and finishable within one academic session. & Correctness argues for hand-authored content and expert-reviewed ontology; the schedule argues for automation. \\
    Supervisor / examiner & Evidence of engineering depth, rigour, and honest evaluation. & An honest evaluation reports unvalidated components, which conflicts with the team's interest in presenting a finished product. \\
    Content rights holders & Their material not redistributed. & The system needs that material as source for retrieval and generation. \\
    \bottomrule
  \end{tabularx}
\end{table}

\subsection{Functional requirements}

\begin{table}[htbp]
  \centering
  \caption{Functional requirements}
  \label{tab:functional}
  \begin{tabularx}{\textwidth}{l>{\hsize=1.6\hsize}Y>{\hsize=0.4\hsize}Y}
    \toprule
    ID & Requirement & Objective \\
    \midrule
    FR1 & A learner can create an account and sign in. & O7 \\
    FR2 & The system presents a catalogue of courses, sections and topics derived from the skill ontology. & O1, O7 \\
    FR3 & A learner can start a diagnostic for a section; the system serves questions one at a time, adaptively selected. & O4 \\
    FR4 & Each response updates the learner's mastery estimate for the tested skill using Bayesian inference conditioned on the question's Bloom level. & O2 \\
    FR5 & A correct response propagates upward, raising ancestor skills to at least the descendant's posterior. & O3 \\
    FR6 & A skill's learning rate is gated on all of its direct prerequisites being mastered. & O3 \\
    FR7 & Mastery views are locked until the section diagnostic is complete. & O4 \\
    FR8 & After the diagnostic, a learner can practise a chosen topic until every skill in it reaches the mastery threshold. & O5 \\
    FR9 & Repeated failure attributable to a missing prerequisite triggers injection of prerequisite-focused questions. & O5 \\
    FR10 & A learner can view their mastery as both a graph and a table, with subject and topic labels. & O7 \\
    FR11 & A learner can reset a section and retake its diagnostic. & O7 \\
    FR12 & Questions render Bangla text with embedded \LaTeX{} mathematics correctly. & O6, O7 \\
    FR13 & The question bank is generated from the source corpus and aligned to skill, Bloom level and topic. & O6 \\
    FR14 & The interface is available in English and Bangla. & O7 \\
    \bottomrule
  \end{tabularx}
\end{table}

\subsection{Non-functional requirements}

\begin{table}[htbp]
  \centering
  \caption{Non-functional requirements}
  \label{tab:nonfunctional}
  \begin{tabularx}{\textwidth}{l>{\hsize=0.55\hsize}Y>{\hsize=1.45\hsize}Y}
    \toprule
    ID & Attribute & Requirement \\
    \midrule
    NFR1 & Correctness & The prerequisite graph must be acyclic; mastery estimates must never place an ancestor below its descendant; these must be checked automatically, not assumed. \\
    NFR2 & Responsiveness & Next-question selection must complete fast enough to feel immediate in a browser session; the search must terminate even when no ideal item exists. \\
    NFR3 & Robustness & Absence of a matching question must degrade gracefully through a defined fallback ladder rather than failing the session. \\
    NFR4 & Maintainability & The ontology must be compiled from a single source of truth by a reproducible build, not hand-edited in its served form. \\
    NFR5 & Verifiability & The engine must be testable end to end without database credentials. \\
    NFR6 & Portability & The system must run on a standard developer machine with no GPU. \\
    NFR7 & Localisation & All learner-facing content must support Bangla, including mathematics. \\
    NFR8 & Security & Secrets must never enter version control; the database service key must remain server-side. \\
    \bottomrule
  \end{tabularx}
\end{table}

\subsection{Constraints}

\textbf{Time:} one academic session, with the team also carrying a full course
load. \textbf{Budget:} effectively zero; every external service had to have a
usable free tier (\Cref{sec:budget}). \textbf{Hardware:} development on
personal laptops with no GPU --- this constraint alone eliminated the original
local-model content-generation approach. \textbf{Data:} no existing learner
response data was available, so no model parameter could be fitted empirically;
all parameters had to be set from pedagogical priors. \textbf{Platform:} web,
to avoid the cost of maintaining native clients. \textbf{Deployment:} a single
process sufficient for demonstration, which is why session state is held in
memory (\Cref{sec:deviations-session}).

\section{Regulatory Requirements and Standards}
\label{sec:standards}

\begin{table}[htbp]
  \centering
  \caption{Standards and regulatory requirements, and how the design responds}
  \label{tab:standards}
  \begin{tabularx}{\textwidth}{>{\hsize=0.6\hsize}Y>{\hsize=0.75\hsize}Y>{\hsize=1.65\hsize}Y}
    \toprule
    Instrument & Why it applies & How the design complies \\
    \midrule
    Data-protection obligations for personal data of minors and young adults
    \cite{bdda2023} & The system stores per-learner performance data; many users
    are 17--19 years old. & Data minimisation: the schema stores only a username
    and per-skill mastery --- no email, phone, address, date of birth or payment
    data. No third-party analytics or advertising trackers are embedded. Deletion
    is supported at section granularity by the retake endpoint. \emph{Gap:} the
    current prototype has no password authentication, so this minimal data is
    not access-controlled; see \Cref{sec:deviations-auth}. \\
    OWASP Top 10 \cite{owasp2021} & The system is a public-facing web
    application with a database behind it. & All database access goes through a
    parameterised client library rather than string-concatenated SQL, addressing
    injection. The privileged service key is held server-side in an environment
    file that is git-ignored and never shipped to the browser, addressing
    sensitive-data exposure. \emph{Gap:} broken access control --- username-only
    login and wide-open CORS are known violations retained deliberately for
    demonstration scope, and are the first items in the production hardening
    list (\Cref{sec:future}). \\
    WCAG 2.1 \cite{wcag21} & The interface is the sole delivery channel and
    excludes users if inaccessible. & Semantic HTML controls, keyboard-operable
    question options, text-based mathematics rendered as accessible markup rather
    than images, and a bilingual interface reducing language exclusion.
    \emph{Gap:} no formal audit against AA criteria was performed
    (\Cref{sec:equity}). \\
    Copyright in source textbooks and question banks & Content is derived from
    copyrighted printed material. & Source material is used as a private
    retrieval corpus, is not redistributed, and is excluded from the public
    repository's distributable artefacts; generated items are new text grounded
    in that corpus rather than reproductions of it. \\
    IEEE/ACM codes of ethics \cite{ieee2020code,acm2018code} & The system makes
    inferences that steer a student's study time during a high-stakes year. &
    Honest reporting of the system's unvalidated components, and refusal to
    claim score improvement without evidence, follow directly from these codes
    (\Cref{sec:codes}). \\
    \bottomrule
  \end{tabularx}
\end{table}

\section{Formulation of Solutions}
\label{sec:alternatives}

Three genuinely different architectures could each satisfy most of the
functional requirements. They differ in where the intelligence lives.

\textbf{Alternative A --- Item Response Theory over a flat item pool.} Model each
\emph{question} with a difficulty parameter and each learner with a single
latent ability parameter; select the next item to maximise information about
that ability. This is the classical computerised-adaptive-testing design.

\textbf{Alternative B --- Deep Knowledge Tracing.} Model the learner's entire
response history with a recurrent network that predicts the probability of
answering the next item correctly \cite{piech2015deep}; select items to
maximise learning gain under the network's predictions.

\textbf{Alternative C --- Bayesian Knowledge Tracing over a prerequisite DAG
(selected).} Model each skill separately with an interpretable four-parameter
hidden-state model, link skills by prerequisite edges, and propagate evidence
along those edges.

\begin{table}[htbp]
  \centering
  \caption{Evaluation of candidate solutions against criteria drawn from the requirements}
  \label{tab:alternatives}
  \begin{tabularx}{\textwidth}{>{\hsize=0.9\hsize}Y>{\hsize=0.7\hsize}Y>{\hsize=0.7\hsize}Y>{\hsize=0.7\hsize}Y}
    \toprule
    Criterion (source) & A: IRT & B: DKT & C: BKT+DAG \\
    \midrule
    Produces an actionable per-skill diagnosis (FR4, and the core problem statement) & No --- one ability scalar cannot say \emph{which} skill is missing & No --- hidden vector is not interpretable per skill & \textbf{Yes} \\
    Can express prerequisite structure (FR5, FR6) & No & Only implicitly, if data suffices & \textbf{Yes, explicitly} \\
    Trainable with zero historical learner data (project constraint) & Needs calibration data per item & Needs large response datasets & \textbf{Yes --- parameters set from pedagogical priors} \\
    Interpretable to a teacher (teacher stakeholder) & Partly & No & \textbf{Yes} \\
    Computational cost on commodity hardware (NFR6) & Low & High (training) & \textbf{Low} \\
    Robust when the question bank is incomplete (project reality) & Poor --- needs calibrated items & Poor --- needs coverage & \textbf{Fair --- fallback ladder plus propagation} \\
    Implementation effort within one session & Medium & High & Medium \\
    \bottomrule
  \end{tabularx}
\end{table}

\noindent\textbf{Conclusion.} Alternative C was selected. The decisive criteria
are the first and third rows. The project's entire premise is telling a student
\emph{which} foundation is missing, which A and B structurally cannot do: A
compresses knowledge to one number, and B's state is a vector with no per-skill
reading. Independently, both A and B require historical response data to
calibrate, and the project began with none and would have none until it had
users --- a circular dependency that C avoids because its four parameters have
defensible pedagogical defaults. B remains the most credible successor once
real response data accumulates, and the engine was deliberately built so that
the model can be replaced without touching its call sites
(\Cref{sec:bkt-model}).
